\chapter{Backupkonzept}
Während dieser Arbeit wird die 3-2-1-1-0-Methode verwendet. Diese ist eine moderne Erweiterung der klassischen 3-2-1-Methode.
In der 3-2-1-Methode stehen die einzelnen Zahlen für folgende Regeln.

\begin{itemize}
    \item [3.] Mindestens drei Kopien
    \item [2.] Auf mindestens zwei verschiedenen Medientypen. (z.B. Lokale Festplatte und Cloud)
    \item [1.] Und mindestens eine Kopie muss an einem geografisch anderen Standort sein.  
\end{itemize}

Das Ziel dieser Methode ist es, einerseits eine sichere Datenenerhältlichkeit zu gewährleisten und ebenfalls bei Problemen keine Daten zu verlieren.

In der modernen 3-2-1-1-0-Methode werden zwei neue Regeln hinzugefügt.

\begin{itemize}
    \item [1.] Mindestens eine Kopie sollte offline sein, also nicht mit dem Netzwerk verbunden und unveränderlich. 
    \item [0.] Die Kopien dürfen keine Fehler aufweisen und werden dafür regelmässig getestet.   
\end{itemize}

Mit den neuen zwei Regeln werden die Daten nun auch gegen Ransomeware oder menschliche Fehler geschützt.
Auch das Testen der Backups wird gefordert, um ungewünschte Überraschungen zu vermeiden.
